%%%%%%%%%%%%%%%%%%%%%%%%%%%%%%%%%%%%
%%%%  Author: Aakash Ravichandran %%%%
%%%%  Description: Introduction to flutter, motivation, and thesis objectives
%%%%%%%%%%%%%%%%%%%%%%%%%%%%%%%%%%%%

\documentclass[../main]{subfiles}
\begin{document}

\chapter{Introduction}
\label{cha:introduction}

\section{Motivation}

In recent decades infrastructure including buildings and bridges has developed rapidly due to urbanisation and the need to 
meet people's requirements. To satisfy this demand, construction structures have been built on a larger scale. This results 
in structures with greater exposure to wind loads. For example, when considering the design of a building, as the height 
increases, the importance of accurate analysis of the wind loads acting on the structure also increases. The same applies to 
the design of most structures such as bridges and wind turbines. \\

In the design of long-span bridges, cable-stayed bridges and suspension bridges are the solutions that satisfy this requirement. 
However, due to their highly optimised nature, these structures become more flexible under human-induced limit states. 
Consequently, they are more susceptible to unsteady wind loads, especially to aeroelastic phenomena such as flutter. 
Flutter is an aeroelastic phenomenon that can cause instability or failure of structures through interaction with wind. 
Until 1940, there was limited focus on this domain, primarily because of an insufficient understanding of wind-structure 
interaction \cite{gimsing2012}. The failure of the Tacoma Narrows Bridge (span 853 m) in the USA in 1940 was a pivotal 
moment in this field. Within four months of its opening, the bridge failed due to flutter at a wind speed of 68 km/h \cite{abbas2017, billah1991}. \\

The key challenge in analysing this flutter behaviour is quantifying the aerodynamic forces that act on the structure. 
In 1934, Theodorsen, a Norwegian American aerodynamicist, developed an analytical solution for aerodynamic forces for 
unsteady flutter analysis of an airfoil \cite{TheodorsenGarrick1934}. This was reformulated by Scanlan in 1971 \cite{scanlan1971} 
with the motivation of applying it to bluff bodies such as bridge decks. This reformulation introduced new shape constants 
known as flutter derivatives. These are a set of functions that quantify the unsteady aerodynamic forces acting on a 
structure in relation to its vibration frequency. \\

CFD simulations or wind tunnel tests are performed to obtain these constants. One method to obtain them with CFD simulation 
is forced motion simulation, where a 2D bridge deck is oscillated at a specific frequency, and the flutter derivatives are 
obtained based on the aerodynamic forces. Since these constants are functions, multiple CFD simulations by varying the vibration 
frequency are required to obtain the flutter derivatives, which is a time-consuming process. \\

Currently, with bridge spans exceeding 2000 m, faster results are needed so that efficient designs can be obtained. 
To overcome the time consumption of the forced motion simulation method, the multi-frequency forced motion vibration method 
was introduced by Lin Huang et al. in 2011 \cite{huang2011}. In this method, instead of applying a single motion frequency, 
one can superpose multiple motions of different frequencies and apply them. Then, using mathematical tools such as the 
Fourier transform, the individual components can be obtained. In this way, the method can theoretically be \(N\) times as 
efficient (where \(N\) is the number of combinations). However, according to the authors, the multi-frequency method reduced 
the simulation time by only 18\%. This was because when using the multi-frequency vibration method, the accuracy of flutter 
derivatives was affected, and therefore more simulation time was required, which resulted in a decrease in the method's 
efficiency. \\

Subsequently, for several years no research focused on this topic. In 2023, G. Martínez-López investigated the multi-frequency 
method to obtain flutter derivatives via curve fitting with LES turbulence modelling \cite{Martínez-López_2023}. The same result 
as with Lin Huang was observed; that is, to obtain accurate flutter derivatives using the multi-frequency method, the simulation 
time must be increased. An important observation from that work is that the selection of the frequency played a key role in the 
method's accuracy. This raises several open issues. First, the reasons why the multi-frequency method fails to achieve its 
theoretical efficiency improvement remain unclear. Understanding the source of this inaccuracy is essential for refining the 
method. Second, the influence of frequency choice on accuracy and the criteria for efficient frequency selection have not been 
established. Furthermore, the extent to which frequency selection affects simulation time and the relevance of any resulting 
inaccuracy for the critical velocity require investigation. Third, additional parameters beyond frequency, such as motion 
amplitude or the number of frequency combinations, may also affect accuracy. Identifying these parameters would contribute 
to a more comprehensive understanding of the method's behaviour.

\section{Objective}
The investigation of the open questions outlined above forms the core objective of this master's thesis. The work follows 
a two-step approach. If inaccuracies appear in the multi-frequency method, they may originate either from the CFD simulations 
or from the mathematical tools used in post-processing. To examine the latter, signal analysis using FFT was performed with 
the SciPy FFT module in Python. To assess potential CFD-related errors, single-frequency simulations were performed to validate 
the principle of superposition on flutter derivatives. Based on the signal analysis, guidelines are developed for selecting 
suitable frequencies. Similarly, based on the single-frequency simulations, guidelines for choosing appropriate motion 
amplitudes are derived. Using these guidelines, two sets of multi-frequency CFD simulations were carried out to evaluate 
the effectiveness of the proposed guidelines. 

\section{Structure of the Thesis}
This document is structured as follows:

\begin{itemize}
    \item \textbf{Chapter 2: Flutter} \\
    This chapter provides the theoretical foundation for understanding flutter analysis. It covers structural wind engineering concepts, 
    the mathematical formulation of flutter using flutter derivatives according to Scanlan's theory, and the identification of aerodynamic forces 
    through forced-motion CFD simulations. Then, the calculation of critical wind speed using aerodynamic and structural properties is presented.
    Finally, the chapter introduces the multi-frequency method for flutter derivative identification.

    \item \textbf{Chapter 3: Signal Analysis} \\
    This chapter investigates the Discrete Fourier Transform (DFT) as a mathematical tool for signal decomposition for multi-frequency method. 
    The analysis examines single-frequency and multi-frequency signal extraction, identifies potential sources of error such as spectral leakage 
    and noise, and provides guidelines for choosing frequencies in the multi-frequency simulations. 

    \item \textbf{Chapter 4: Single Frequency Simulations} \\
    This chapter conducts a study on single-frequency simulations. It includes a mesh convergence study to determine optimal mesh sizing, 
    followed by a linearity assessment examining the influence of motion amplitude on flutter derivatives. The chapter analyzes aerodynamic forces, 
    spectral characteristics, and pressure distributions to establish amplitude thresholds for linear behavior.

    \item \textbf{Chapter 5: Multi-Frequency Simulations} \\
    This chapter applies the findings from signal analysis and single frequency simulations to conduct multi-frequency CFD campaigns. 
    Two sets of multi-frequency simulations are performed using the formulated frequency and amplitude criteria. The results are compared 
    with single-frequency analysis to validate the proposed guidelines and demonstrate efficiency improvements.

    \item \textbf{Chapter 6: Conclusion and Outlook} \\
    The chapter concludes with recommendations for 
    future research to extend the guidelines to additional cross-sections and turbulence models.

\end{itemize}

\end{document}
